\documentclass[slideColor]{prosper}
\usepackage{amsmath,amsfonts,amsthm}
%\usepackage{epsfig}
\include{Qcircuit}


\newcommand{\E}{\mathbf{E}}
%\newcommand{\R}{\mathbb{R}}

\newcommand{\cA}{\mathcal A}
\newcommand{\cB}{\mathcal B}
\newcommand{\cH}{\mathcal H}
\newcommand{\C}{\mathbb C}
\newcommand{\TVD}{\mathrm{TVD}}

\renewcommand{\>}{\rangle}
\newcommand{\<}{\langle}


%%%%%%%%%%%%%%%%%%%%%%%%%%%%%%%%%%%%%%%%%%%%%%%%%%%%%%%%%%%%

\title{Thermalizing Quantum Systems \&\\ Evaluating Partition Functions\\ with a Quantum Computer}
\author{Pawel M. Wocjan\\[1ex]School of Electrical Engineering and Computer Science\\ University of Central Florida\\Orlando}
\email{wocjan@eecs.ucf.edu}

\begin{document}

\maketitle
%%%%%%%%%%%%%%%%%%%%%%%%%%%%%%%%%%%%%%%%%%%%%%%%%%%%%%%%%%%%%%%%%%%%%%%

\begin{slide}{Joint work with}
\begin{itemize}
\item David Poulin (University of Sherbrooke) 

\medskip
arXiv preprint 0905.2199

\bigskip

\item Chen-Fu Chiang (University of Central Florida)

\bigskip\bigskip
\item this work has been supported by the National Science Foundation
grants CCF-0726771 and CCF-0746600
\end{itemize}
\end{slide}

%%%%%%%%%%%%%%%%%%%%%%%%%%%%%%%%%%%%%%%%%%%%%%%%%%%%%%%%%%%%%%%%%%%%%%%%%

\begin{slide}{Definitions}
\begin{itemize}
\vspace{-0.25cm}
\item let $H$ be a Hamiltonian with spectral decomposition
\[
H = \sum_{a=1}^D E_a |\psi_a\>\<\psi_a|\,, \quad E_a \in [0,1]
\]

\item the partition function $Z_\beta$ at inverse temperature $\beta$ is
\[
Z_\beta := \sum_{a=1}^D e^{-\beta E_a}
\]

\item the thermal state $\rho_\beta$ at temperature $\beta$ is
\[
\rho_\beta := \sum_{a=1}^D \frac{e^{-\beta E_a}}{Z_\beta} |\psi_a\>\<\psi_a|
\]

\medskip
\item 
\end{itemize}

\end{slide}

%%%%%%%%%%%%%%%%%%%%%%%%%%%%%%%%%%%%%%%%%%%%%%%%%%%%%%%%%%%%%%%%%%%%%%%%%

\begin{slide}{Tasks}
\begin{itemize}
\item \textcolor{blue}{Evaluating parition functions}

\medskip\noindent
given $H$, $\beta$, and $\epsilon$, output a random value $\tilde{Z}_\beta$ such that
\[
\Pr \Big( \big| \tilde{Z}_\beta - Z_\beta \big| < \epsilon Z_\beta \big) > \frac{3}{4}
\]

\bigskip
\item \textcolor{blue}{Thermalizing quantum systems}

\medskip\noindent
given $H$, $\beta$, and $\epsilon$, prepare a quantum state $\tilde{\rho}_\beta$ with
\[
\frac{1}{2} \| \rho_\beta - \tilde{\rho}_\beta \|_{\mathrm{tr}} \le \epsilon 
\]

\bigskip
\item
\end{itemize}
\end{slide}

%%%%%%%%%%%%%%%%%%%%%%%%%%%%%%%%%%%%%%%%%%%%%%%%%%%%%%%%%%%%%%%%%%%%%%%%%

\begin{slide}{Thermalizing Quantum Systems}
\begin{itemize}
\item given $H$, $\beta$, and $\epsilon$, prepare a quantum state $\tilde{\rho}_\beta$ with
\[
\frac{1}{2} \| \rho_\beta - \tilde{\rho}_\beta \|_{\mathrm{tr}} \le \epsilon \quad \mbox{(*)}
\]

\item the above task was studied by Terhal and Divincenzo in Phys. Rev. A 61, 022301

\begin{itemize}
\item constructed a TCP whose fix-point is the thermal state $\rho_\beta$ in an idealized
setting (perturbation theory)

\medskip
\item 
did not prove how the number of iterations has to be chosen to obtain a state $\tilde{\rho}_\beta$ satisfying (*)

\medskip
\item numerical simulations suggests that this algorithm may outperform classical algorithms
\end{itemize}
\end{itemize}
\end{slide}

%%%%%%%%%%%%%%%%%%%%%%%%%%%%%%%%%%%%%%%%%%%%%%%%%%%%%%%%%%%%%%%%%%%%%%%%%

\begin{slide}{Our Quantum Algorithm}
\begin{itemize}
\item given $H$, $\beta$, and $\epsilon$, our quantum algorithm makes it possible to prepare $\tilde{\rho}_\beta$ with
\[
\frac{1}{2} \| \rho_\beta - \tilde{\rho}_\beta \|_{\mathrm{tr}} \le \epsilon 
\]

\bigskip
\item its running time is
\[
\sqrt{\frac{D}{Z_\beta}} \, \cdot \, poly(\log D, 1/\epsilon, \beta)
\]
\end{itemize}
\end{slide}

%%%%%%%%%%%%%%%%%%%%%%%%%%%%%%%%%%%%%%%%%%%%%%%%%%%%%%%%%%%%%%%%%%%%%%%%%

\begin{slide}{Classical Hamiltonians}
\begin{itemize}
\item let $\Omega$ be a set whose elements $x$ correspond the states of some classical system

\medskip
\item let $E : \Omega\rightarrow [0,1]$ denote the energy function, assigning to each state $x$ its energy $E(x)$

\medskip
\item let 
\[
H = \sum_{x\in\Omega} E_x |x\>\<x|
\]

\medskip 
\item thermalizing such classical systems is easier because $H$ is diagonal in the computational basis $\{ |x\> \}_{x\in\Omega}$
\end{itemize}

\end{slide}

%%%%%%%%%%%%%%%%%%%%%%%%%%%%%%%%%%%%%%%%%%%%%%%%%%%%%%%%%%%%%%%%%%%%%%%%%

\begin{slide}{Metropolis Algorithm}
\begin{itemize}
\item given such classical Hamiltonian $H$, we can always construct a Markov chain $P$ whose limiting distribution is equal to $\rho_\beta$

\medskip
\item the spectral gap $\delta$ determines how fast we can approach the Boltzmann distribution

\medskip
\item the run time scales like $1/\delta$

\medskip
\item often, it is not possible to determine (or bound) the spectral gap $\Rightarrow$ heuristics

\end{itemize}
\end{slide}

%%%%%%%%%%%%%%%%%%%%%%%%%%%%%%%%%%%%%%%%%%%%%%%%%%%%%%%%%%%%%%%%%%%%%%%%%

\begin{slide}{Quantized Metropolis Algorithm}
\begin{itemize}
\item given such classical Hamiltonian $H$, there is a ``quantized'' Metropolis algorithm preparing $|\tilde{\varphi}_\beta\>$ with
\[
\big\| |\varphi_\beta\> - |\tilde{\varphi}_\beta\> \big\| \le \epsilon
\]
where $|\varphi_\beta\>$ is a ``quantized'' version of the Boltzmann distribution
\[
|\varphi_\beta\> := \sum_{x\in \Omega} \sqrt{\frac{e^{-\beta E_x}}{Z_\beta}} |x\>
\]

\item the run times scales like $\sqrt{1/\delta}$

\medskip
\item 
proved in Somma, Boixo, Barnum, and Knill, PRL 101, 130504 and W. and Abeyesinghe, PRA 78, 032311
\end{itemize}
\end{slide}

%%%%%%%%%%%%%%%%%%%%%%%%%%%%%%%%%%%%%%%%%%%%%%%%%%%%%%%%%%%%%%%%%%%%%%%%%

\begin{slide}{Grover's State Engineering}
\begin{itemize}
\item prepare the state 
\[
\frac{1}{\sqrt{D}}\sum_{x\in\Omega} |x\> \otimes |E_x\> \otimes |0\>
\]

\bigskip
\item apply a controlled rotation to obtain the state
\[
\frac{1}{\sqrt{D}}\sum_{x\in\Omega} |x\> \otimes |E_x\> \otimes \Big( \sqrt{e^{-\beta E_x}} |0\> + \sqrt{1-e^{-\beta E_x}} |1\> \Big)
\]

\end{itemize}
\end{slide}

%%%%%%%%%%%%%%%%%%%%%%%%%%%%%%%%%%%%%%%%%%%%%%%%%%%%%%%%%%%%%%%%%%%%%%%%%

\begin{slide}{Grover's State Engineering}
\begin{itemize}

\item apply Grover to increase the overlap with the $|0\>$ component and measure in the basis $|0\>$, $|1\>$ 


\medskip
\item if we obtain $|0\>$, then the post-measurement state is
\[
\sum_{x\in\Omega} \sqrt{\frac{e^{-\beta E_x}}{Z_\beta}} |x\> \otimes |E_x\> \otimes  |0\> 
\]

\medskip
\item the expected run time to obtain this state scales like
\[
\sqrt{\frac{D}{Z_\beta}}
\]

\end{itemize}
\end{slide}

%%%%%%%%%%%%%%%%%%%%%%%%%%%%%%%%%%%%%%%%%%%%%%%%%%%%%%%%%%%%%%%%%%%%%%%%%

\begin{slide}{Our Quantum Algorithm}
\begin{itemize}
\item is a generalization of the algorithm based on Grover's state engineering to \textcolor{red}{quantum} Hamiltonians

\bigskip
\item let $H$ be an arbitrary Hamiltonian with spectral decomposition
\[
H = \sum_{a=1}^D E_a |\psi_a\>\<\psi_a|\,, \quad E_a \in [0,1]
\]

\end{itemize}
\end{slide}

%%%%%%%%%%%%%%%%%%%%%%%%%%%%%%%%%%%%%%%%%%%%%%%%%%%%%%%%%%%%%%%%%%%%%%%%%

\begin{slide}{Error Sources}
\begin{itemize}
\item given $H$, $\beta$, and $\epsilon$, our quantum algorithm makes prepares $\tilde{\rho}_\beta$
\[
\frac{1}{2} \| \rho_\beta - \tilde{\rho}_\beta \|_{\mathrm{tr}} \le \epsilon 
\]

\medskip
\item the errors come from three sources:
\[
\| \rho                 - \tilde{\rho}         \|_{\mathrm{tr}} \le 
\| \rho                 - \rho_{\textrm{sim}}  \|_{\mathrm{tr}} +
\| \rho_{\textrm{sim}}  - \rho_{\textrm{prec}} \|_{\mathrm{tr}} +
\| \rho_{\textrm{prec}} - \rho_{\textrm{fail}} \|_{\mathrm{tr}} +
\| \rho_{\textrm{fail}} - \tilde{\rho}          \|_{\mathrm{tr}} 
\]

\medskip
\begin{itemize}
\item $\rho_{\textrm{sim}}$ imperfect simulation of $\exp(i H)$

\medskip
\item $\rho_{\textrm{prec}}$ limited precision of phase estimation

\medskip
\item $\rho_{\textrm{fail}}$ nonzero failure probability of phase estimation
\end{itemize}
\end{itemize}
\end{slide}

%%%%%%%%%%%%%%%%%%%%%%%%%%%%%%%%%%%%%%%%%%%%%%%%%%%%%%%%%%%%%%%%%%%%%%%%%

\begin{slide}{Imperfect Hamiltonian Simulation}
\begin{itemize}
\item assume that $H$ is sparse and let $U = \exp(i H)$

\medskip
\item using techniques for simulating Hamiltonian time evolutions, we can implement $U_{\textrm{sim}}$ with
\[
\| U - U_{\textrm{sim}} \| \le \epsilon_{\textrm{sim}}
\]

\medskip
\item the run time scales like 
\[
poly(log D, 1/\epsilon_{\textrm{sim}}, s)
\]
\end{itemize}
\end{slide}

%%%%%%%%%%%%%%%%%%%%%%%%%%%%%%%%%%%%%%%%%%%%%%%%%%%%%%%%%%%%%%%%%%%%%%%%%

\begin{slide}{Technical Arguments}
\begin{itemize}
\vspace{-0.25cm}
\item matrix logarithm is Lipschitz continuous:
\[
\| U - U_{\textrm{sim}} \| \le \epsilon_{\textrm{sim}}
\]
\quad implies that there is a Hamiltonian $H_{\textrm{sim}}$ with
\[
\| H - H_{\textrm{sim}} \| \le O(\epsilon_{\textrm{sim}})
\]

\item thermal states are H\"older continuous:
\[
\| H - H_{\textrm{sim}} \| \le O(\epsilon_{\textrm{sim}})
\]
\quad implies that the corresponding thermal states satisfy
\[
\| \rho - \rho_{\textrm{sim}} \|_{\textrm{tr}} \le O(\sqrt{\beta \epsilon_{\textrm{sim}}})
\]

\item we assume $U_\textrm{sim}=U$ (and keep track of the error)
\end{itemize}
\end{slide}

%%%%%%%%%%%%%%%%%%%%%%%%%%%%%%%%%%%%%%%%%%%%%%%%%%%%%%%%%%%%%%%%%%%%%%%%%

\begin{slide}{Maximally Entangled State}

\begin{itemize}
\item let
\[
|\Phi\> := \frac{1}{\sqrt{D}} \sum_{a=1}^D |i\> \otimes |i\>
\]

\item let
\[
V = \sum_{a=1}^D |\psi_a\>\<a|
\]

\item it is important that we can express $|\Phi\>$ as
\[
|\Phi\> = \sum_{a=1}^D |\psi_a\> \otimes |\psi^T_a\>
\]

this follows from $(V\otimes I_D)|\Phi\>=(I_D\otimes V^T)|\Phi\>$

\medskip
\item 
\end{itemize}

\end{slide}

%%%%%%%%%%%%%%%%%%%%%%%%%%%%%%%%%%%%%%%%%%%%%%%%%%%%%%%%%%%%%%%%%%%%%%%%%

\begin{slide}{Phase Estimation with Boosting}

\begin{itemize}
\item recall that $U=\exp(iH)$

\bigskip
\item $\Rightarrow$ we have
\begin{eqnarray*}
(U \otimes I_D) |\psi_a\> \otimes |\psi^T_a\> & = & e^{i E_a} |\psi_a\> \otimes |\psi^T_a\> \\
&& \\
(U \otimes I_D) |\Phi\>                       & = & \sum_{a=1}^D e^{i E_a} |\psi_a\> \otimes |\psi^T_a\>
\end{eqnarray*}

\item ${}$ $\Rightarrow$ the energy levels $E_a$ of $H$ are mapped bijectively

${}$ \phantom{$\Rightarrow$} onto the phases $e^{i E_a}$ of $U$

\medskip
\item $\Rightarrow$ apply phase estimation to estimate the energy levels

\end{itemize}
\end{slide}

%%%%%%%%%%%%%%%%%%%%%%%%%%%%%%%%%%%%%%%%%%%%%%%%%%%%%%%%%%%%%%%%%%%%%%%%%

\begin{slide}{Phase Estimation with Boosting}

\begin{itemize}
\medskip
\item use phase estimation with $O(1/\epsilon_{\textrm{prec}})$ ancillas and 

run $O(\log(1/\epsilon_{\textrm{fail}}))$ boosting rounds

\medskip
\item $\Rightarrow$ this transforms $|\Phi\>$ into the state
\[
\sum_{a=1}^D |\psi_a\> \otimes |\psi^T_a\> \otimes \Big( c_a^- |E_a^-\> + c_a^+ |E_a^+\> \Big) + |\textrm{junk}\>
\]
where
\begin{itemize}
\item $E_a^+ - E_a  \le \epsilon_{\textrm{prec}}$, \quad $E_a - E_a^- \le \epsilon_{\textrm{prec}}$

\medskip
\item $|c_a^+|^2 + |c_a^-|^2=1$

\medskip
\item $\| |\textrm{junk}\> \| \le \epsilon_{\textrm{fail}}$
\end{itemize}
\end{itemize}
\end{slide}

%%%%%%%%%%%%%%%%%%%%%%%%%%%%%%%%%%%%%%%%%%%%%%%%%%%%%%%%%%%%%%%%%%%%%%%%%

\end{document}



\begin{slide}{FPRAS}
\begin{itemize}
\item the problem of computing the partition function exactly is too difficult; it's \#P-hard

\medskip
\item $\Rightarrow$ consider fully polynomial randomized approximation schemes

\medskip
\item a FRAPS

\begin{itemize}
\item outputs a random number $\hat{Z}$ satisfying
\[
\Pr\Big( (1-\epsilon)Z(\beta) \le \hat{Z} \le (1+\epsilon) Z(\beta) \Big) \ge 3/4
\]
where $\epsilon\in (0,1)$ determines the desired accuracy

\medskip
\item runs in time polynomial in $\log(|\Omega|)$ and $1/\epsilon$
\end{itemize}
\end{itemize}
\end{slide}

%%%%%%%%%%%%%%%%%%%%%%%%%%%%%%%%%%%%%%%%%%%%%%%%%%%%%%%%%%%%%%%%%%%%%%%%%

\begin{slide}{Simulated Annealing}
\begin{itemize}
\item choose a cooling schedule $\beta_0 < \beta_1 < \cdots < \beta_{\ell+1}$ with $\beta_0=0$ and $\beta_{\ell+1}=\beta$

\medskip
\item express the desired quantity as a telescoping product
\[
Z(\beta) = \frac{Z(\beta_{\ell+1})}{Z(\beta_\ell)} \cdot \frac{Z(\beta_\ell)}{Z(\beta_{\ell-1})} \cdots \frac{Z(\beta_1)}{Z(\beta_0)} \cdot Z(\beta_0) 
\]

\medskip
\item observe that $Z(\beta_0)$ is trivial since $Z(\beta_0)=|\Omega|$

\medskip
\item $\Rightarrow$ estimate the ratios $\alpha_i:=Z(\beta_{i+1})/Z(\beta_i)$ 
\end{itemize}
\end{slide}

%%%%%%%%%%%%%%%%%%%%%%%%%%%%%%%%%%%%%%%%%%%%%%%%%%%%%%%%%%%%%%%%%%%%%%%%%

\begin{slide}{Estimation via Boltzmann Sampling I}
\begin{itemize}
\item denote by $\pi_i=\big(\pi_i(x) : x\in\Omega\big)$ the Boltzmann distribution at inverse temperature $\beta_i$
\[
\pi_i(x) = \frac{e^{-\beta_i E(x)}}{Z(\beta_i)}
\]

\item assume we can sample from $\pi_i$

\medskip
\item ensure that each ratio $\alpha_i:=Z(\beta_{i+1})/Z(\beta_i)$ is bounded from below by a constant, say $1/2$

\medskip
\item $\Rightarrow$ it suffices to take a small number of samples $\sim \pi_i$ to estimate $\alpha_i$

\end{itemize}
\end{slide}

%%%%%%%%%%%%%%%%%%%%%%%%%%%%%%%%%%%%%%%%%%%%%%%%%%%%%%%%%%%%%%%%%%%%%%%%%

\begin{slide}{Estimation via Boltzmann Sampling II}
\begin{itemize}
\item let $X_i \sim \pi_i$

\medskip
\item let $\Delta\beta_i = \beta_{i+1} - \beta_i$ 

\medskip
\item define a new random variable 
\[
Y_i = e^{-\Delta\beta_i E(X_i)}
\]
 
\item the expected value $\E(Y_i)$ is equal to
\[
\sum_{x \in \Omega}
\pi_i(x) 
e^{-\Delta\beta_i E(x)}  =
\sum_{x \in \Omega}
\frac{e^{-\beta_i E(x)}}{Z(\beta_i)} 
e^{(-\beta_{i+1} + \beta_i) E(x)}  
= 
\alpha_i
\]
\item $\Rightarrow$ $Y_i$ is an unbiased estimator for $\alpha_i$
\end{itemize}
\end{slide}

%%%%%%%%%%%%%%%%%%%%%%%%%%%%%%%%%%%%%%%%%%%%%%%%%%%%%%%%%%%%%%%%%%%%%%%%%

\begin{slide}{Estimation via Boltzmann Sampling III}
\begin{itemize}
\item draw $O(l/\epsilon^2)$ samples of $X_i$ and compute the mean $\bar{Y}_i$

\medskip
\item $\Rightarrow$ the random variable $\bar{Y}=\bar{Y}_0 \bar{Y}_1 \cdots \bar{Y}_{\ell}$ satisfies
\[
\Pr\Big(
(1-\epsilon) \alpha < \bar{Y} < (1+\epsilon) \alpha
\Big)
\ge 3/4
\]
where $\alpha=\alpha_0\alpha_1\cdots\alpha_\ell$

\medskip
\item $\Rightarrow$ the total number of samples is
\[
O\Big(\ell^2 / {\red \epsilon^2} \Big)
\] 
\end{itemize}
\end{slide}

%%%%%%%%%%%%%%%%%%%%%%%%%%%%%%%%%%%%%%%%%%%%%%%%%%%%%%%%%%%%%%%%%%%%%%%%%

\begin{slide}{Sampling with Markov Chains}
\begin{itemize}
\item in general, we are not able to sample directly from $\pi_i$ 

\medskip
\item $\Rightarrow$ construct a Markov chain $P_i$ such that 

\begin{itemize}
\item its stationary distribution is equal to $\pi_i$ 

\medskip
\item its spectral gap $\delta$ is large

\begin{itemize}
\item $\Rightarrow$ we have to simulate  
\[
\tilde{O}(1/ {\red \delta})
\] 
steps of $P_i$ to obtain one sample from $\tilde{\pi}_i$

\medskip
\item $\tilde{\pi}_i$ has to be $O(\epsilon^2/\ell^2)$ close to $\pi_i$
\end{itemize}

\end{itemize}

\end{itemize}
\end{slide}

%%%%%%%%%%%%%%%%%%%%%%%%%%%%%%%%%%%%%%%%%%%%%%%%%%%%%%%%%%%%%%%%%%%%%%%%%

\begin{slide}{Quantum Speed-Up}
\begin{itemize}
\item classical complexity
\[
\tilde{O}\Big(
\ell^2 /\big(\delta \epsilon^2\big) 
\Big)
\]

\item quantum complexity
\[
\tilde{O}\Big(
\ell^2/\big(\sqrt{\delta} \epsilon\big)
\Big)
\]

\item $1/\delta \rightarrow 1\sqrt{\delta}$ is due to the quadratic relation between the spectral gaps of $P_i$ and the phase gaps of the corresponding quantum walks $W(P_i)$ 

\medskip
\item $\epsilon^2 \rightarrow \epsilon$ is due to quantum estimation of expected values of observables
\end{itemize}
\end{slide}

%%%%%%%%%%%%%%%%%%%%%%%%%%%%%%%%%%%%%%%%%%%%%%%%%%%%%%%%%%%%%%%%%%%%%%%%%

\begin{slide}{Classical Walk}
\begin{itemize}
\item $P_i$ is a time-reversible Markov chain 
\[
P_i=\Big(p_i(x,y)\Big)_{x,y\in\Omega}
\]
with stationary distribution 
\[
\pi_i=\Big(\pi_i(x)\Big)_{x\in\Omega}
\]

\item $\pi_i$ is the unique left eigenvector of $P_i$ with eigenvalue $1$
\[
\pi_i P_i = \pi_i
\]

\end{itemize}
\end{slide}

%%%%%%%%%%%%%%%%%%%%%%%%%%%%%%%%%%%%%%%%%%%%%%%%%%%%%%%%%%%%%%%%%%%%%%%%%

\begin{slide}{Spectral Gap}

\begin{itemize}

\item let $1=\lambda_0 > \lambda_1 > \ldots > \lambda_{|\Omega|-1}$ be the spectrum of $P_i$

\bigskip
\item the spectral gap 
\[
\delta := 1 - \lambda_1
\]

\medskip
determines how fast $\pi_i$ is approached from any initial state
\end{itemize}
\end{slide}

%%%%%%%%%%%%%%%%%%%%%%%%%%%%%%%%%%%%%%%%%%%%%%%%%%%%%%%%%%%%%%%%%%%%%%%%%

\begin{slide}{Quantum Update}

\begin{itemize}
\item a quantum update for $P_i$ is any unitary $U_i$ acting on $\C^{|\Omega|}\otimes\C^{|\Omega|}$ with
\begin{equation*}
U_i \Big( |x\> \otimes |0\> \Big) = 
|x\> \otimes 
\sum_{y\in\Omega} \sqrt{p_i(x,y)} |y\> 
\end{equation*}
for some fixed state $0\in\Omega$

\bigskip
\item classical complexity: the number of times we apply $P_i$

\bigskip
\item quantum complexity: the number of times we apply $U_i$ or $U_i^\dagger$
\end{itemize}
\end{slide}

%%%%%%%%%%%%%%%%%%%%%%%%%%%%%%%%%%%%%%%%%%%%%%%%%%%%%%%%%%%%%%%%%%%%%%%%%

\begin{slide}{Quantum Walk}
\begin{itemize}
\item the quantum walk $W(P_i)$ is realized by applying $U_i$ twice and $U_i^\dagger$ twice (and also some other simple unitaries)

\bigskip
\item the unique eigenvector of $W(P_i)$ with eigenvalue $1$ is
\[
|\pi_i\> = \sum_{x\in\Omega} \sqrt{\pi_i(x)} |x\>
\]

\item
$|\pi_i\>$ is a coherent version of the limiting distribution $\pi_i$
\end{itemize}
\end{slide}

%%%%%%%%%%%%%%%%%%%%%%%%%%%%%%%%%%%%%%%%%%%%%%%%%%%%%%%%%%%%%%%%%%%%%%%%%

\begin{slide}{Quadratic Relation}
\begin{itemize}
\item the phase gap $\Delta$ of $W(P_i)$ is
\[
\min\{ |\varphi| : \mbox{ $e^{i\varphi}$ is an eigenvalue of $W(P_i)$, $e^{i\varphi}\neq 1$} \} 
\]

\medskip
\item the quadratic relation between the phase and spectral gaps
\[
\Delta \ge \sqrt{\delta}
\]
is at the heart of quantum speed-ups of many search problems
\end{itemize}
\end{slide}

%%%%%%%%%%%%%%%%%%%%%%%%%%%%%%%%%%%%%%%%%%%%%%%%%%%%%%%%%%%%%%%%%%%%%%%%%

\begin{slide}{Primitives}
\begin{itemize}
\item let
\[
\Pi_i := |\pi_i\>\<\pi_i| \quad \mbox{and} \quad \Pi_i^\perp = I - \Pi_i
\]

\medskip
\item we can implement approximate versions of 
\begin{eqnarray*}
R_i         & = & - \Pi_i             + \Pi_i^\perp \\
S_i         & = & e^{-2\pi j/3} \Pi_i + \Pi_i^\perp \quad \mbox{and} \quad S_i^\dagger 
\end{eqnarray*}

\medskip
by applying the quantum walk $W(P_i)$
\[
\tilde{O}\Big(
\frac{1}{\sqrt{\delta}}
\Big)
\]
times
\end{itemize}
\end{slide}

%%%%%%%%%%%%%%%%%%%%%%%%%%%%%%%%%%%%%%%%%%%%%%%%%%%%%%%%%%%%%%%%%%%%%%%%%

\begin{slide}{Quantum Estimation of the Ratios I}
\begin{itemize}
\item let $A_i$ be the observable
\[
A_i = \sum_{x\in\Omega} e^{-\Delta\beta_i E(x)} |x\>\<x|
\]

\medskip
\item we have 
\[
\<\pi_i|A_i|\pi_i\> = \alpha_i
\]
\end{itemize}
\end{slide}

%%%%%%%%%%%%%%%%%%%%%%%%%%%%%%%%%%%%%%%%%%%%%%%%%%%%%%%%%%%%%%%%%%%%%%%%%

\begin{slide}{Quantum Estimation of the Ratios II}
\begin{itemize}
\item let 
\[
V_i := \sum_{x\in\Omega} |x\>\<x| \otimes
\left(
\begin{array}{cc}
 \sqrt{e^{-\Delta\beta_i E(x)}}   & \sqrt{1-e^{-\Delta\beta_i E(x)}} \\
-\sqrt{1-e^{-\Delta\beta_i E(x)}} & \sqrt{e^{-\Delta\beta_i E(x)}}
\end{array}
\right)
\]

\medskip
\item let 
\[
|\psi_i\> = V_i \Big( |\pi_i\> \otimes |0\> \Big) \quad \mbox{ and } \quad \Gamma := I \otimes |0\>\<0|
\]

\medskip
\item $\Rightarrow$ we have
\begin{eqnarray*}
\<\psi_i|\Gamma|\psi_i\> & = & \<\pi_i|A_i|\pi_i\> \\
& = & \alpha_i
\end{eqnarray*}
\end{itemize}
\end{slide}

%%%%%%%%%%%%%%%%%%%%%%%%%%%%%%%%%%%%%%%%%%%%%%%%%%%%%%%%%%%%%%%%%%%%%%%%%

\begin{slide}{Quantum Estimation of the Ratios III}
\begin{itemize}
\item we apply quantum phase estimation to 
\[
G = \big(2|\psi_i\>\<\psi_i)-I\big) \big( 2\Gamma - I \big) \quad\mbox{and}\quad
|\psi_i\> = |\pi_i\>|0\>
\]

\item $\Rightarrow$ we get a random variable $Q_i$ such that
\[
\hspace{-0.5cm}
\Pr\left( \Big(1-\frac{\epsilon}{2(\ell+1)}\Big)\alpha_i \le Q_i \le \Big(1+\frac{\epsilon}{2(\ell+1)}\Big)\alpha_i \right) \ge 1 - \frac{1}{16(\ell+1)}
\]

\medskip
\item it suffices to apply the reflection $R_i$ around $|\pi_i\>$
\[
O\Big(
\frac{\ell}{\epsilon}
\Big)
\]
times
\end{itemize}
\end{slide}

%%%%%%%%%%%%%%%%%%%%%%%%%%%%%%%%%%%%%%%%%%%%%%%%%%%%%%%%%%%%%%%%%%%%%%%%%

\begin{slide}{Quantum Estimation of the Ratios IV}
\begin{itemize}
\item let $Q$ be the random variable 
\[
Q=Q_0 Q_1 \cdots Q_\ell
\]

\item $\Rightarrow$ we have
\[
\Pr\Big( (1-\epsilon)\alpha \le Q \le (1+\epsilon)\alpha \Big) \ge \frac{15}{16}
\]

\item it suffices to apply operators from $\{W(P_0),\ldots,W(P_\ell)\}$
\[
O\Big(
\frac{\ell^2}{\sqrt{\delta}\epsilon}
\Big)
\]
times
\end{itemize}
\end{slide}

%%%%%%%%%%%%%%%%%%%%%%%%%%%%%%%%%%%%%%%%%%%%%%%%%%%%%%%%%%%%%%%%%%%%%%%%%

\begin{slide}{Preparation of Quantum Samples I}
\begin{itemize}
\item the fact that the ratios $\alpha_i$ are bounded from below by $1/2$ implies 
\[
|\<\pi_i|\pi_{i+1}\>|^2 \ge \frac{1}{2}
\]

\medskip
\item $\Rightarrow$ we can drive $|\pi_i\>$ to $|\tilde{\pi}_i\>$ by applying operators from the set $\{S_i,S_i^\dagger,S_{i+1},S_{i+1}^\dagger\}$
\[
\tilde{O}(1)
\] 
times

\medskip
\item this is based on Grover's $\frac{\pi}{3}$-fixed point search

\end{itemize}
\end{slide}

%%%%%%%%%%%%%%%%%%%%%%%%%%%%%%%%%%%%%%%%%%%%%%%%%%%%%%%%%%%%%%%%%%%%%%%%%

\begin{slide}{Preparation of Quantum Samples II}
\begin{itemize}
\item starting from $|\pi_0\>$ (easy because uniform superposition), we can prepare any $|\tilde{\pi}_i\>$ by
\[
|\pi_0\> \rightarrow |\tilde{\pi}_1\> \rightarrow \cdots \rightarrow |\tilde{\pi}_i\> \rightarrow \cdots \rightarrow |\tilde{\pi}_\ell\>
\]

\medskip
\item $\Rightarrow$ it suffices to apply the operators from the set $\{S_k,S_k^\dagger : k = 0,\ldots,i\}$ $\tilde{O}(i)$ times

\medskip
\item $\Rightarrow$ the complexity of approximating
\[
|\pi_0\> \otimes |\pi_1\> \otimes \cdots \otimes |\pi_\ell\>
\]
is 
\[
\tilde{O}\Big(
\ell^2
\Big)
\]
\end{itemize}
\end{slide}

%%%%%%%%%%%%%%%%%%%%%%%%%%%%%%%%%%%%%%%%%%%%%%%%%%%%%%%%%%%%%%%%%%%%%%%%%

\begin{slide}{Summary}

\hspace{5cm}
\Qcircuit @C=0.9em @R=0.9em {
\lstick{|\pi_0\>}                                                                                  & \gate{\tilde{\alpha}_0}      & \qw \\ 
\lstick{|\pi_0\> \rightarrow |\tilde{\pi}_{1}\>}                                                   & \gate{\tilde{\alpha}_1}      & \qw \\ 
                                                                                                   &                   &     \\
\lstick{|\pi_0\> \rightarrow |\tilde{\pi}_{1}\> \rightarrow \cdots \rightarrow |\tilde{\pi}_{l}\>} & \gate{\tilde{\alpha}_{\ell}} & \qw    
}

\bigskip
for each $i=0,\ldots,\ell$,
\begin{itemize}
\item the complexity of preparing an approximate version $|\tilde{\pi}_i\>$ of $|\pi_i\>$ is $\tilde{O}(\ell)$

\medskip
\item the complexity of estimating $\alpha_i$ is 
$\tilde{O}\Big(
\frac{\ell}{\sqrt{\delta} \epsilon}
\Big)$ 
\end{itemize}

$\Rightarrow$ overall complexity $\tilde{O}\Big(
\frac{\ell^2}{\sqrt{\delta} \epsilon}
\Big)$ 

\end{slide}

%%%%%%%%%%%%%%%%%%%%%%%%%%%%%%%%%%%%%%%%%%%%%%%%%%%%%%%%%%%%%%%%%%%%%%%%%

\begin{slide}{Future Research}

\begin{itemize}
\item quantum speed-up of estimating permanents of matrices with non-negative entries

\medskip
\item quantum speed-up of estimting the volume of a convex polytope?

\medskip 
\item quantum speed-up of other classical approximation algorithms?

\medskip
\item estimating the partition functions of quantum Hamiltonians???
\end{itemize}

\end{slide}

%%%%%%%%%%%%%%%%%%%%%%%%%%%%%%%%%%%%%%%%%%%%%%%%%%%%%%%%%%%%%%%%%%%%%%%%%

\end{document}